
%\section[Congruences, Graphs, Reachability \& Distinguishability]{Relating Congruences, Orbital Graphs, Subset Reachability and Distinguishability}


%\subsection{Sufficient Conditions} % for Complete Reachability and Sync-Maximality}

%\subsection{Permutation Groups, Graphs and \texorpdfstring{$G$}{G}-congruences} %%and Block Systems} %statt block system congruences

%\subsection{Permutation Groups, Graphs and Congruences} 


\todo{das als eigene section und als untersection davon primitive gruppe? dann subclasses ganz am schluss}

\todo{auch im wesentlichen von mir, und jordan verallgemeinert.}

\todo{diese jordan satz verallgemeinerung diskutieren.}

The main result of this section is Theorem~\ref{sync_max:thm:no_congruence_equivalences}
stating various equivalences relating  the property 
that two points cannot
be congruent for a non-trivial $G$-congruence in a given permutation group $G$
(or that they are never both contained in a non-trivial block)
to complete reachability, distinguishability
of subsets and generating sets for the full symmetric group.\todo{oder detecting if a symmetric group, wie in isaacs
paper über generating sets beschrieben.}

\todo{begriff orbits nutzen und definieren?}




\begin{toappendix}
Before stating the main result, the section begins with two lemmata giving sufficient
conditions for complete reachablity and distinguishability of subsets. The proofs
rely heavily on constructing graphs on which a given permutation group acts as
an automorphism group.







% pivotal results? crucial results

\todo{https://link.springer.com/content/pdf/10.1007/s10801-016-0703-9.pdf
lineare schranke für konstant abbildungen.}

\todo{im lemma definition versteckt}

\todo{circular simple idempotent binär klassifzieren, äquiavlenz kürzetses syncwort auch, aber folgt auch durch einfache transformation
auf Cerny automaten}

\todo{semigroup statt monoid}
\begin{lemma}
\label{sync_max:lem:sufficient_compl_reach}
Let $M \le \mathcal T_n$ be a transformation semigroup. \todo{erwähnen compl reachable semigroup ist monoid}
Consider the directed graph $\Gamma_{\text{SI}}(M)$ with vertex set $[n]$ \todo{im text davor definieren.}
and edge set
\[ 
\{ (a,b) \mid a \ne b \text{ and there exists a simple idempotent in $M$ mapping $a$ to $b$ } \}.
\]
If the graph is strongly connected, then $M$
is completely reachable
and every two distinct $2$-sets are distinguishable.
\end{lemma}
\begin{proof}
The defined graph is a subgraph of $\Gamma_1(M)$.
Hence if it is strongly connected, complete reachability
is implied with Proposition~\ref{todo}. 
Now, the graph has the following property:
if there exists a path $x_1, x_2, \ldots, x_n$
in the graph then the mapping that results by composing the simple idempotents
that induce the edges along the path maps $[n]$ to $[n] \setminus \{x_1, x_2, \ldots, x_{n-1}\}$.\todo{beweis}
Let $\{a,b\}, \{c,d\}$ be two distinct $2$-sets
and consider a shortest path $x_1, x_2, \cdots, x_n$
with $x_1 = a$ and $x_n = b$.
As we have choosen a shortest path, it contains no loops
and so $|\{x_1, x_2, \ldots, x_n\}| = n$.
Then the corresponding element maps $\{a,b\}$ to $\{b\}$.\todo{genauer:
	$[n]$ auf $[n] \setminus \{ a, x_2, \ldots, x_{n-1} \}$
	damit folgt nicht $\{a,b\}$ auf $\{b\}$
	eher mit $\{x_1, \ldots,x_n\}$ auf $\{x_n\}$ (gleich da rest self-loops hier),
	vgl. blatt 9.7 (1) bleistiftskizze
}
If $\{c,d\} \cap \{x_1, x_2, \ldots, x_n\} = \emptyset$,
then $\{c,d\}$ is mapped to $\{c,d\}$.
If there exists $i \in \{1,\ldots,n\}$
such that $c = x_i$ but $d \notin \{x_1, x_2, \ldots, x_n\}$,
then $\{c,d\}$ gets mapped to $\{b,d\}$, similarly
if $d = x_i$ and $c \notin \{x_1, x_2, \ldots, x_n\}$.
If there exists $1 \le i < j \le n$
such that $x_i = c$ and $x_j = d$ (or vice versa), \todo{nochmal genauer hinschreiben.}
then the mapping composed from
the path $x_i, x_{i+1}, \ldots, x_j$ maps $\{c,d\}$ to $\{d\}$
but as $\{a,b\}\cap \{x_i,x_{i+1},\ldots,x_j\} = \emptyset$
maps $\{a,b\}$ to itself.
\end{proof}
\end{toappendix}

\todo{group of units "controls" some properties}

\todo{mit kernel von $f$ ausdrücken, gemeinsame verfeinerung von knogruenz auf $G$
und dem kernel, falls das nur equality dann die bedingungen erfüllt?}



\begin{comment} 
\begin{lemma}
\label{sync_max:lem:no_congruence_ab}
Let $f : [n] \to [n]$ be a function of defect one %rank $n - 1$
and $a,b \in [n]$ be distinct such that $\applyfunc{a}{f} = \applyfunc{b}{f}$.
% muss nicht transitiv sein, da sonst ja eine g-kongruenz bewahrt
Let $G \le S_n$ be a permutation group
admitting no non-trivial $G$-congruence for which $a$ and $b$ 
are congruent (or equivalently no non-trivial block containing $a$ and $b$) and such that there exists $g \in G$
with $\applyfunc{a}{g} = b$.
Then
\begin{enumerate}
	\item the transformation monoid $\langle G \cup \{f\} \rangle$ is completely reachable,
	
	\todo{sagen distinghished, und nur distinct hier gesagt!}
	\item for any two distinct $2$-sets $\{a,b\}, \{c,d\}$ there exists a transformation
	in $\langle G \cup \{f\} \rangle$
	mapping precisely one to a singleton set but not the other.
	
\end{enumerate}
\end{lemma}
\begin{proof} \todo[inline]{annahme $f(b) = b$ noch weglassen, da man potenz so wählen kann , dass auf $b$ landet. braucht nur, dass bild permutiert. wenn 
	
	$f(a) = f(b) = c$ und $f^n(c) = b$ (da permutiert und $b$ sei im bild),
	dann $f^{n+1}(a) = f^{n+1}(b) = f^n(c) = b$ 
	
	und wenn bild auf identität, dann nur $f^i(a) = f^i(c) = c$.
	(hier $i = n + 1$ möglich)
}
By Lemma~\ref{sync_words:lem:idempotent} we can assume $f$ is a simple idempotent mapping.\todo{oder idempotent of defect one immer sagen?}
Set $M = \langle G \cup \{f\}\rangle_M$.
The mapping $f$ induces the edge $(a,b)$ in $I(M)$.
By Equation~\ref{rank-defect-inqeuliaty} we can deduce that $G$ is precisely the group of units of $M$.
As $G$ is a subgroup of the automorphism group of $I(M)$ (Proposition~\ref{sync_max:prop:G_automorhism_on_graphs})
the connected components (without taking the direction of the edges into account)
of $I(M)$ form congruence classes for $G$.
As $(a,b)$ is an edge in $I(M)$,
both are contained in the same connected component.
Then by assumption all elements must be congruent, i.e., $I(M)$ must be connected.
With Lemma~\ref{sync_max:lem:ab_non_congruent_then_two_orbits}
and as there exists an element mapping $a$ to $b$,
the group $G$ is transitive.
Hence $I(M)$ is a vertex-transitive graph.
By Theorem~\ref{todo} this implies that $I(M)$ is strongly connected.
Then the claim follows with Theorem~\ref{sync_max:thm:IS_SCC}.
\end{proof}
\end{comment}


\todo{wie verhält sich permutation wenn mit zyklus über alle elemente verknüpft, bilder dazu, nutzen
in circular chapter, und für transposition}

\todo{colon nutzen.}

\todo{die fussnote über $n > 2$ hier dran. und das lemma im einführungstext erwähnen.}

\todo{btw das für "nur $a,b$" point stabilizer nicht maximal ist klar,
da transitiv gleich operation auf punktstabs und
ja durchaus blöcke erlaubt, und die entsprechen halt obergruppen und denn 
nebenklassen in der obergruppe. wäre ja auch komisch wenn alle punktstabs, da transitiv
als konjugierte von einem, maximal, weil dann wäre es transitiv.

das vielleicht noch schreiben im primitiv teil?

damit auch nennung dieses begründen.


beispielgruppe

C4, dann 1 und 2 nie kongruent, punktstab von 1, identität, nicht maximal
}

\begin{comment}
\todo[inline]{FAALSSSSSCCHCCCCCC}
\begin{lemma}
\label{sync_max:lem:ab_non_congruent_then_transitive}
Let $G \le \mathcal S_n$ with $n > 2$ and $a,b \in [n]$ be distinct.
If there does not exist a non-trivial block containing $a$
and $b$, then $G$ acts transitively on $[n]$.
\end{lemma}
\begin{proof}
Suppose $G$ does not act transitively.
Set $\Gamma_a = \{ \applyfunc{a}{g} \mid g \in G \}$
and $\Gamma_b = \{ \applyfunc{b}{g} \mid g \in G \}$.
If $\Gamma_a \cap \Gamma_b = \emptyset$, 
then $\{ a \} \cup \Gamma_b$ is a non-trivial block
containing~$a$ and~$b$. \todo{genauer?}
Otherwise, if $\Gamma_a = \Gamma_b$, then, as $G$ 
is not transitive by assumption, $[n]\setminus \Gamma_a \ne \emptyset$ 
and $\Gamma_a$ is a non-trivial block containing $a$ and $b$.

\todo[inline]{
	
	Scheint falsch, $\Gamma_b \cup \{a^g\}$ ist doch möglich mit $a^g \ne a$,
	aber $\Gamma_a \cup \Gamma_b$ ist ein Block, also alles.
	Damit ist jedes Element entweder  im Orbit von $a$
	oder $b$!!!
}
\end{proof}
\end{comment} 






\todo{etwas mehr zu $n \ge 5$ sagen. eigentlich nur wegen der vorletzten bedingung mit den
disjunkten mengen, ein remark welche äquvilenzen auch für klienere n gelten.}

\todo{subset fragen, aus neumann/morgenstern und literatur dazu angeben.}


\todo{fussnote erzeuger monoid oder semigroup egal bei gruppen} 


\todo[inline]{

Siehe DLT submission, dort für SI-automaten gezeigt (mit beliebig vielen simple idempotents) wenn G transitiv.

Allgemeines prinzip dahinter in perm.-gruppensprache:
a1,b1, …, ak,bk, wenn alle kongruent/dann universelle kongruenz, also alle kongruent (ist hinter der aussage für SI-automaten).	
}


\begin{theorem} \todo{transitive brauche ich nicht, und kleine fälle $n \ge 4$, gibt es da sonderfälle?}
\label{sync_max:thm:no_congruence_equivalences} \todo{ja okay, aber warum $n \in \{3,4\}$ ausgeschlossen?} 
Let $G \le_G \mathcal S_n$ with\footnote{For $n \le 2$ there cannot exist a non-trivial
congruence on $[n]$ and some equivalences fail to hold true.}
$n \ge 5$ \todo{wenn nicht transitiv, fall $a$ und $b$ in verschiedenen bahnen vno $G$.
aber das widerspricht nicht der aussage!

zwei primzyklen, $f$ einen auf den anderen. dann nicht completely reachable da nicht SCC,
und nicht sync-max!
aber in dem fall volle kongreunz auf einen und triviale auf anderen!
}
and $a,b \in [n]$ be distinct. Suppose there exists $h \in G$ such that $\applyfunc{a}{h} = b$.
Then the following are equivalent:

\todo{einige der vorderen äquiavlenzen brauche dieses $g$ nicht, dass auch erwähnen in satz danach.}

\begin{enumerate}[label=(\arabic*)]




%\item For every transformation semigroup homomorphism $\varphi : [n] \to [m]$
% we have $\applyfunc{a}{\varphi} \ne \applyfunc{b}{\varphi}$.

\item \label{sync_max:thm:ab_completely_reachable} For every [idempotent] $f \colon [n] \to [n]$ of defect one with $f(a) = f(b)$,
the transformation semigroup $\langle G \cup \{f\} \rangle_S$ is completely reachable.

\item \label{sync_max:thm:ab_subset-pair} For every [idempotent] $f \colon [n] \to [n]$ of defect one with $f(a) = f(b)$,
the transformation semigroup $\langle G \cup \{f\} \rangle_S$ is completely distinguishable.


\item \label{sync_max:thm:ab_subsets_different_cardinality} For every [idempotent] $f\colon [n] \to [n]$ of defect one with $f(a) = f(b)$
and distinct non-singleton subsets $A, B \subseteq [n]$ 
there exists $t \in \langle G \cup \{f\} \rangle_S$
with $|\applyfunc{A}{t}| \ne |\applyfunc{B}{t}|$. 




\todo{transformation statt function sagen}

\item \label{sync_max:thm:ab_distinguishes_2_sets} For every [idempotent] $f\colon [n]\to [n]$ of defect one with $f(a) = f(b)$,
in the transformation semigroup $\langle G \cup \{f\} \rangle_S$
pairs of distinct $2$-sets are distinguishable.

%\item \label{sync_max:thm:ab_distinguishes_all_subsets} For every [idempotent] $f\colon [n] \to [n]$ of defect one with $f(a) = f(b)$,
% the transformation monoid $\langle G \cup \{f\} \rangle$ distinguishes %all non-empty and 
% non-singleton
% subsets.\todo{die ganzen anderen, disjunkte mengen etc?}

\item \label{sync_max:thm:ab_distinguishes_disjoint_2_sets} For every [idempotent] $f\colon  [n] \to [n]$ of defect one with $f(a) = f(b)$,
in the transformation semigroup $\langle G \cup \{f\} \rangle_S$ 
pairs of disjoint $2$-sets are distinguishable.


\todo{erzeugte halbgruppe oder monoid?}

\item \label{sync_max:thm:2-hom} For every [idempotent] $f\colon [n] \to [n]$ of defect one with $f(a) = f(b)$
and for every $c,d \in [n]$ there exists $t \in \langle G \cup \{f\} \rangle_S$ 
such that $\applyfunc{\{c,d\}}{t} = \{a,b\}$. \todo{auch für completely subset-pair wenn erzeugermenge mit genau einer nicht-permuation.}  


\todo{ \item There exists a function $f : [n] \to [n]$ of defect one with $f(a) = f(b)$
	
	Existenzaussagen!
}

\todo{usw für alle bedingungen. mit der bedingung kann man beliebige funktion
	nehmen und testen bzw dann auch für alle nicht der fall
	
	klar, gibt nur zwei simple idempotente, und ziel sit egal
	
	
	Damit folgt einfacher Algrithmus wenn nur permutation und genau eine vom defekt one.
	halt orbital von (a b) ausrechnen und gucken ob connected. wenn nicht connected da für alle-aussage.
	
	
	der algorithmus vielleicht effizienter?
	
	
	Let $f$ be given. $G \cup \{f\}$ compleltey reachable iff orbital von a und b connected in $G$.
	wenn orbital connected, dann compl reachalbe, da für alle funktionen
	wenn compl reachable mit f, dann existiert eine funktion (exist aussage), also connected
	
	existiert aussage auch klar, da konkrete funktion eigentlich irrelevant und nur die punkte a und b
	entscheidend
	
	
	
	für existenzaussage beweis: gleich, aber in der konstruktion der funktion, nimm simple idempotetn
	aus $f$ und nutze, dass $G$ auf kongruenz als äquiavlenz operirt, also ist auch beliebige funktion
	äquivalenz-treu auf kongruenz; nie das potenzieren von $fg$ macht dieses argument kaputt, also fordern
	äquivalenz-treue funktion
	
	-> also Existenz-Verstärkung nur für simple idempotent (oder äquiavlenz-treue bzgl g-congruence, aber das als
	bemerkung); neue abbildung kann ja auch auf orbital operieren; aber beachte man kann annehmen, dass
	er bild permutiert, weil dieses g dahinterscahlten, dann äquivalenzt-treu genau dann wenn ausgangsabbildung
	dies war, konjugation mit defekt one abbildung? beie "fasern" betrachten und neues tupel hinzufügen.
}

\item \label{sync_max:thm:ab_constant_map} For every [idempotent] $f\colon [n] \to [n]$ of defect one with $f(a) = f(b)$,
the transformation semigroup $\langle G \cup \{f\} \rangle_S$ contains a constant map.



\todo{maximal subgroup definieren.}


\todo{definition point stabilizer nur im theorem machen und vorne wegnehmen.}

% ne, das gilt wohl nicht.
%\item \label{sync_max:thm:point_stabs} The point stabilizer subgroups $G_a$ and $G_b$ are maximal subgroups in $G$.


\todo{is synchronizing, ne den begriff nicht verwenden, da irreführend
	mit synchroniing group.}

\end{enumerate}

\end{theorem}
\begin{proof} First, we show the equivalences $\ref{sync_max:thm:ab_no_nontrivial_congruence} 
\Leftrightarrow
\ref{sync_max:thm:ab_blocks} 
\Leftrightarrow 
\ref{sync_max:thm:undirected_orbital}
\Leftrightarrow
\ref{sync_max:thm:ab_orbitals} 
\Leftrightarrow 
\ref{sync_max:thm:partition}$.
Then we show the following implications.

\begin{center} 
% https://tikzcd.yichuanshen.de/#N4Igdg9gJgpgziAXAbVABwnAlgFyxMJZABgBoBGAXVJADcBDAGwFcYkQAzLjkAX1PSZc+QigBMFanSat23bnwEgM2PASIBmUsSkMWbRJ3k9+g1SM2kxumQaMLeUmFADm8IqA4AnCAFskZCA4EEjkppw+-ojkNMFIYuHefqGxIYgajrxAA
\begin{tikzcd}[arrows=Rightarrow]
	&                                           \ref{sync_max:thm:ab_subset-pair} \arrow[r] & \ref{sync_max:thm:ab_subsets_different_cardinality} \arrow[r] & \ref{sync_max:thm:ab_distinguishes_2_sets} \arrow[r] & \ref{sync_max:thm:ab_distinguishes_disjoint_2_sets} \arrow[r] 
	& \ref{sync_max:thm:2-hom}  \arrow[rd] &  \\
	\ref{sync_max:thm:ab_orbitals} \arrow[ur] \arrow[rrr] & & & \ref{sync_max:thm:ab_completely_reachable} \arrow[rrr] & & &  \ref{sync_max:thm:ab_constant_map} \arrow[r] & \ref{sync_max:thm:ab_no_nontrivial_congruence}  
\end{tikzcd}
\end{center}

\todo{in diskussion erwähnen egienschaft strongly factorizable semigroup}

\todo{satz sollte auch für $n = 4$ gelten? und wenn keine disjunkten $2$-sets gilt es auch?}


\todo{mit dem lemma primitive dann transitive falls mehr als zwei punkt, nimm drei $Gamma_a \cup \Gamma_b = \Gamma_a \cup \Gamma_c = \Gamma_b \cup \Gamma_c$ implies alle gelich $[n]$}



\todo{anderes argument
	% im primitiven fall
	% folge g_1(a) \in B und g_2(b) in A, dann nochmal anwenden g_2 usw
	% b durch g_1g_2 ... immer in A irgendwann identität, dann wäre b \in A, transitiv nutzen b auf a, muss a nach B da primitiv
	
	nicht primitiv, jeweils urbild unter g1 g2 usw
	
	amit aus 1 und 2.
}

\todo{äquivalent oder congruenct sagen?}

\todo{sagen hier unterscheidung monoid/semigroup nicht whcihtig, da immer monoid}

\todo{index $S$ oder $M$ dranschreiben bzw. sagen on purpose weggelassen.}


\begin{comment} 
% mit den funktionensätzen zeigen? als eine simple idempotent gleich
% direkt mit graph und connected-lemma
%
% f überall durch transposition ersetzen, dann hat man praktisch stehts zwei mengen A, B, A währe bild in G \cup f, B ausgeschlossen

\todo{$f$ statt $f_i$!}
Set $t = (a ~ b)$.
Let $S_0 = \applyfunc{[n]}{(g_0 f_1 g_1 \cdots g_{n-1}f_n g_n)}$ % schreibe h = g_0 f_1 ... 
and \todo{$g_n$ zuerst,linksaction}
\[ % es muss |S_0| = |S_1| = ... usw sein
S_i = \begin{cases}
	\applyfunc{( S_{i-1} \setminus \{b\} ) \cup \{a\}}{g_i^{-1}} & \text{if } a \in \applyfunc{[n]}{(g_i f_i g_{i+1} \cdots g_{n-1} f_n g_n)} \\ % \{a,b\} \subseteq \applyfunc{[n]}{(g_i f_i g_{i+1} \cdots g_{n-1} f_n g_n)} \\
	\applyfunc{S_{i-1}}{g_i^{-1}} & \text{if } a \notin \applyfunc{[n]}{(g_i f_i g_{i+1} \cdots g_{n-1} f_n g_n)}
\end{cases}
\]
We claim that for each $i \in \{1,\ldots,n\}$ % oder allgemeiner für alle i <= j?
we have $S_0 = \applyfunc{S_i}{(g_0 t g_1 \cdots g_{i-1} t g_i)}$.

Suppose $S_0 = \applyfunc{S_i}{(g_0 t g_1 \cdots g_i t)}$. If 
$a \in \applyfunc{[n]}{(g_{i+1} f_{i+1} g_{i+2} \cdots g_{n-1} f_n g_n)}$,
then $S_{i+1} = ( S_i \setminus \{b\} ) \cup \{a\}$.


\todo{schreibe die verkettung direkt rein, ohne applyfunc, damit ist das makro ja sinnnlos, oder ein makro concat}

%% wie vorher





% h_i definieren, mit h_0 = identität
Suppose $S_0 = \applyfunc{S_i}{(g_0 t g_1 \cdots g_{i})}$.
Set $T_i = \applyfunc{[n]}{(g_{i+1} \cdots g_{n-1} f_n g_n)}$.
If $a \in T_i$, then $S_{i+1} = (\applyfunc{S_i}{g_i^{-1}} \setminus \{g_i^{-1}(b)\}) \cup \{g_i^{-1}(a)\}$.
So
\[
\applyfunc{S_{i+1}}{(g_0 t g_1 \cdots g_{i-1}tg_i)}
= \applyfunc{( S_i \setminus \{b\} ) \cup \{a\}}{(g_0 t g_1 \cdots g_{i-1}t)}
= \applyfunc{ S_i \cup \{b\} }{(g_0 t g_1 \cdots g_{i-1})}
\]





Set $T_i = (g_i f g_{i+1} \cdots g_{n-1} f g_n)([n])$
and $S_0 = g_0^{-1}(T_0)$ and, for $i > 0$,
\[ % t( .. ) hier doch überflüssig?
S_i = \begin{cases}
	t(g_{i-1}^{-1}(S_{i-1}))                               & \text{if } \{a,b\} \cap T_i   = \emptyset \\ 
	(g_{i-1}^{-1}(S_{i-1}) \setminus \{b\} ) \cup \{a\} & \text{if } \{a,b\} \cap T_i \ne \emptyset.
\end{cases}
\]
We claim that $T_0 = (g_0 t g_1 \cdots g_{i-1} t g_{i} t)(S_i)$. For $i = 0$
this is clear. So assume $0 < i < n$. 
If $\{a,b\} \cap T_{i+1} = \emptyset$, then
\[
(g_0 t g_i \cdots g_{i-1} t g_i t )(S_{i+1})
= (g_0 t g_i \cdots g_{i-1} t g_i)(S_{i})
\]
\end{comment}
\end{proof}


\todo{
kongruenezn bei cirular mit untergruppen entsprechnung?
}


\todo[inline]{

dixon mortimer aufgabe

Delta definiere simple idempotent, mit fall im finite, erzeuge idempotente alpha auf beta 


anders folgern: gibt produkt simple idempotent mit bild teilmenge die alpa enthält, aber nicht beta


oder teilmengen separierungseigenschaft

----

f rk n - 1


zwei idempotent mappings daraus, wenn die beiden elemente in keiner kongruenz, dann sync-max & completely reachable (da auch SCC, weil gruppe transitiv) 

muss nichtmal idempotent herstellen, nur das ziel von dem einen


geht doch bei beliebigem f und operation durch konjugation (im prinzip gleichen functionalgraph!)



für das schließen enifach kante mit ursprung exluced, ziele irgendwas im bild, dann kann man auch entlang so schließen.

- aber exluded element wird u.U. bewegt, also kann man nicht entlang des pfades so schließen


wirklich nur "simple idempotent" graph
}

\todo{kann man auch immer aussuchen, welches 2-set auf einerset? diskutieren -> ne, dann 2-homogen, remark,
und beispiel primitive aber nicht 2-hom gruppe C_p nennen.}

\todo{distinguishability 
impliziert höchstens 2 orbits.

ne, auch transitiv, als remark?

}

\todo{pair kreterium, schreiben alle $2$-sets statt pairs.}



\todo[inline]{
Dieser Remark bezog sich noch auf älteren Beweis vor dem 27.11. Ist der Kommentar noch gültig?
\begin{remark}
That Condition~\ref{sync_max:thm:ab_distinguishes_2_sets}
implies Condition~\ref{sync_max:thm:ab_constant_map}
can be shown
similarly
as the implication ``$\ref{sync_max:thm:ab_distinguishes_disjoint_2_sets} \Rightarrow \ref{sync_max:thm:ab_constant_map}$''
was shown but without using Theorem~\ref{thm:separating_neumann}
in the following way.
First, note that  Condition~\ref{sync_max:thm:ab_distinguishes_2_sets}
implies transitivity.
For assume we have a non-trivial orbit $\Gamma$. 
\todo{orbit definieren}
Let $s \in [n] \setminus \Gamma$ and $t,r \in \Gamma$ with $t \ne r$. \todo{oder aus dem aneren zwei, einer muss ja zwei enthalten.}
Then we can distinguish \todo{genauer}
$\{s,t\}$ and $\{s, r\}$.
But one of those subset can only be mapped to $\{a,b\}$
if $a \in \Gamma$ or $b \in \Gamma$, but $\{a,b\} \not\subseteq \Gamma$, \todo{anders schreibe, schnitt kardinalittät eins.}
as the other element
is from a different orbit. \todo{genauer}
However, then take three elements from $\Gamma$
or from $[n] \setminus \Gamma$, which is possible as $n \ge 5$,
and for two distinct $2$-sets. Then neither of these can be mapped ot $\{a,b\}$.

Now let $\{s,t\} \subseteq [n]$
be any $2$-set. As $G$
is transitive, \todo{dafür brauche ich nichtmal transitiv, nur jeder orbit mindens drei elemente.}
there exists $g \in G$
with $\applyfunc{\{s,t\}}{g} \ne \{s,t\}$.
Now, distinguish both sets and proceed as
in the proof 
``$\ref{sync_max:thm:ab_distinguishes_disjoint_2_sets} \Rightarrow \ref{sync_max:thm:ab_constant_map}$''
above.
\end{remark}
}


\todo[inline]{

Noch ein paar Bemerkungen (27.11.2022)




1) Statt Für-Alle-f-Aussagen, gelten auch Exists-f-Aussagen?

exists f so dass ... statt für alle

f idempotent annehmen mit $f^g$ (also als produkt von konjugieten) 
formel. dann kongruenzklasse um a,b die existiert, dann für die nicht sync-bar, also keine konstante map.
damit folgt, wenn es $f$ gibt, so dass erzeugte Halbgruppe konstante Abbildung, dann $a$ und $b$ nie kongruent?



2) (so in LATA-paper neue version)
f(0) = f(1) = 1
pqr kongruenz auf eine feste zahl 1 (kürzestes erzeuger oder transitiv); oder einfach äquivalent, wenn auf einen mapbar mit anderem draußen
-> damit transformation 2-hom; dann anderer beweis jordan (idemp durch transposition tauschen)
(transitivität der pqr-kongruenz auch mit A.t singleton und ein q.t \not A.t?)
p ~ q <=> forall r \notin {p,q} exists p.f = q.f und r.f != p.f

neuer beweis jordan, zyklus mit nerode-kongruenz, zyklus schrittweise erweitern, damit einen großen zyklus (aus hamilton-pfad)
-> orbital graph hat hamilton pfad, induktiv bauen: auf teilmenge $A$ gibt es ihn, dann auch pfad der rausgeht; ne, der kann aber irgendwo mitten in bestehenden pfad starten


oder irgendeine kongruenz definieren, so dass man alle mengen unterscheiden kann. also dass man auf graphentheoretische argumente verzichten kann.

und so, dass entscheidungsverfahren rausnehmbar begründbar.



3) 
\ref{sync_max:thm:partition} mit nerode und transitiv wie in rystsov-paper dass graph core? remark mit alternativen beweisen


4) 
2-set distinguishable => sync mit den small Satz~\ref{sync_max:prop:2-sets-dist-SCC-implies-sync} 
und da transitiv (alternativer beweis, remark)

5) 
p,q,r bedingung, kongruenz definieren? Festem r definieren p ~ q 
bemerkung für alternativen beweis

6) bedingung |t(S)| != |t(T)| (different cardinality), immer ein idempotent 
gilt immer, wenn erzeugermenge mit nur einem simple idempotent? da dann kürestes produkt wählen.
also in dem fall dass zu completely subset-pair distinguishable äquivalent?

}

\todo{prim group distniguishes pairs(!) schreiben
oder 2-sets pairwise istinguishable
}





\todo[inline]{

diese sclussweise in dem folgenden todo-teil zeigt doch, dass
$\langle G \cup f \rangle$ $2$-homogen operiert, und das characterisiert
primitivität (bzw a und b nicht kongruent). ist doch ein nettes ergebnis.
damit auch, dass man stets von zwei $2$-sets sich aussuchen kann, welches
auf singleton (da sogar auf den 2-er sets injektiv operiert, mit anderen ergebnissen).

weiterhin folgt damit einfach jordan, wenn man f durch transposition austauscht.

{s,t}, {p,q} eins auf singelton, also auf {a,b} mit g.
sei dies ${s,t}$, dann betzarchte ${p,q}^g$ und ${p,q}$,
davon eins auf ${a,b}$, damit in gleicher bahn wie {s,t}

}

\begin{toappendix}

\todo[inline]{

neuer beweis jordan:

transposition (1 2) dazu.
dann 2 2-er set bahnen (mit funktion f(1) = f(2) dazu, oder graph argument direkt und pfad entspricht
mapping von 2-set auf 2-set).
dann andere transpsition dazu die aus beiden tauscht, um 2-er bahnen zu einer 2-er bahn zu machen.
aber ist nicht mehr mit der ausgangstransposition!

{a,s} {b,s} einer auf {a,b}, dann aber auch anderer, da transpostiion die tauscht.


also wenn in anderer bahn, dann kein a oder b drin
{s,t} auf {p,q} in einem trennenden paar, wo anderes in {a,b}-bahn (also {x,y},{s,t} und {x,y} auf {a,b})
dann aber betrachte diese beiden, dann davon einer auf {a,b}, also in der bahn, und damit der andere
-> damit 2-hom.

kann man damit auch 3-hom zeigen?
}


\todo{das folgende besser einordnen, noch jordan damit beweisen

primitive genau dann wenn höchstens zwei $2$-sets orbits.}

\begin{theorem} %gliech auf {a,b}
$G$
is $2$-homogenous if and only if for any distinct $2$-sets $\{s,t\}, \{p, r\}$
there exists one mapping $\{s,t\}$ to a singleton.
\end{theorem}
\begin{proof}
This is equivalent to mapping $\{s,t\}$ to $\{a,b\}$, as $\{s,t\}$
was arbitrary, we have only one orbit on the $2$-sets.

FALSCH, kann ja $f$ mit benutzen. (aber gegen transpostiion austauschen für jordan)
\end{proof} 

\begin{theorem}
$a$ und $b$ never congruent  iff we have at most two orbits on the $2$-sets.
\end{theorem}
\begin{proof}
one orbit klar, da dann $2$-homogenous.

gibt es mehr als zwei orbits, so nimm zwei $2$-sets aus orbits
die nicht zu $\{a,b\}$ gehören


direkt, mit kongruenzen?

$C_4$ hat auch zwei $2$-set bahnen?

ne, hab wohl übersehen dass man $f$ ja auch "dazwischen" benutzen darf,
aber kann man gegen transposition austauschen, dann beweis für jordan.
\end{proof}

\todo[inline]{
prmitiiv, höchstens zwei 2-er set bahnen?

wenn mehr, dann die zwei aus denen nicht zu {a,b} wären nicht unterschiedbar!


äquivalent primitiv genau dann wenn höchstens zwei 2-er set-bahnen für alle a,b?
(ist aussage über orbital graphen)


---
damit jordan neu, transpisiton elment zwischen den 2-er bahnen -> 2-hom -> damit jede transposition auf jede andere.


2-hom {1,2,3...,n}
mit < (n+1 n+2) >
dort kann man alle trennen, und nicht transitiv

{1,n+1}, {2,n+2} nicht

außer für f(1) = f(n+1)

dann die in 1,2,3,4 nicht
}
\end{toappendix}
